\subsection{nml\_sem}
\label{SEC:nml_sem}
\inp*[nmlSEM]{}[nml\_sem][]

This \hyperref[SEC:User_Input]{namelist block}
is used to specify the global controls of the Synthetic Eddy Method.
%
Additional controls for specifying the SEM on individual boundaries
are contained in the \inptext{nml\_bc} namelist block
described in \sref{SEC:nml_bc}.

\begin{nmlblocklst}[
,label={LST:nml_sem}
,caption={Namelist nml\_sem.}
]
&nml_sem
  Use_Improved_SEM_Normalization = T
  Use_Profile_For_Initialization = F
  Use_Random_Seed_From_Restart   = F
  Use_Spherical_Eddy_Volume      = F
  Include_Thermo_Fluctuations    = F
  Box_Num_Eddies_Multiplier      = 1.0
  Use_Restart_Eddies             = T
  Tent_Function_Continuity       = -1
/
\end{nmlblocklst}

%==========================================================================

\begin{nmldescitem}
Use_Improved_SEM_Normalization = T
\end{nmldescitem}
\inp*[nmlSEM::UseImprovedSEMNormalization]{}[Use\_Improved\_SEM\_Normalization][]

\begin{nmldescbody}
This uses an exponential weighted moving average of the eddy concentration as
the normalization factor for the unscaled fluctuating velocity field instead of
the value $\sqrt{\mathscr{V}_{\mathrm{box}} / \left(\mathscr{V}_{\mathrm{eddy}}
N_{\mathrm{eddy}}\right)}$ used by the classic SEM formulation.
%
More details can be found in \linecite{Skillen_2016}.
\end{nmldescbody}

%==========================================================================

\begin{nmldescitem}
Use_Profile_For_Initialization = F
\end{nmldescitem}
\inp*[nmlSEM::UseProfileForInitialization]{}[Use\_Profile\_For\_Initialization][]

\begin{nmldescbody}
Specifies whether the SEM profile will be used to initialize the
\emph{entire volume} of the flow domain based on
the $y$ coordinate value and the profile height variable.
%
This is very experimental and only works for very simple geometries
like a flat plate.

%%Hide this, because few users will need it.
%\NOTE
%This does some more specialized flow initializations if used with the
%THX3 test case.

%\FIXME{
%What are the prerequisites regarding structured or unstructured grids,
% and uniformity in a specific direction?
%}
\end{nmldescbody}

%==========================================================================

\begin{nmldescitem}
Use_Random_Seed_From_Restart = F
\end{nmldescitem}
\inp*[nmlSEM::UseRandomSeedFromRestart]{}[Use\_Random\_Seed\_From\_Restart][]

\begin{nmldescbody}
When using an SEM boundary condition, the SEM data will also be
written to the restart files including the seed used for the random number
generator.

If \inptext{Use\_Random\_Seed\_From\_Restart} is true, the seed for the
random number generator will be read from
the restart file and used to reset the random number generator to exactly
where it left off when the restart file was written, as if the simulation
never stopped.

If it is false, the current date and time will be used to create a new
``random'' seed for the random number generator. This might have a very small
effect on the ``randomness'' of the generated strength and locations of the
synthetic eddies in the SEM BC.
%
However, the total number of time steps needed for a simulation will almost
certainly prevent this from having any kind of measurable effect on the
simulation results.
%
\NOTE Very rarely (but it does happen), the SEM can randomly cause a
simulation to die due to an NaN or other floating point exception.
Setting this input to false is the only way to get past one of these
occasions without modifying the solution with limiting \andor filtering.
\end{nmldescbody}

%==========================================================================

\FINALFIX{\newpage}
\begin{nmldescitem}
Use_Spherical_Eddy_Volume = F
\end{nmldescitem}
\inp*[nmlSEM::UseSphericalEddyVolume]{}[Use\_Spherical\_Eddy\_Volume][]

\begin{nmldescbody}
This input determines if the volume of each synthetic eddy is treated as a
sphere (true) or as a box (false).
%
%It really only has an effect on the number of eddies that will be used
It affects the number of eddies that will be used
for an individual SEM boundary,
%
%[{Number of eddies}] = \inptext{Box\_Num\_Eddies\_Multiplier} $\times$
%                       [{Volume of SEM box}]
%                     / [{Volume of synthetic eddy}].
%
\begin{gather*}
     [\text{Number of eddies}]
  =  [\inptext{Box\_Num\_Eddies\_Multiplier}] \times
     \frac{ [\text{Volume of SEM box}] }
          { [\text{Volume of synthetic eddy}] }
\end{gather*}
%
where the volumes are described in \sref{SEC:SEM:theory:Nomenclature}.

\end{nmldescbody}

%==========================================================================

\begin{nmldescitem}
Include_Thermo_Fluctuations = F
\end{nmldescitem}
\inp*[nmlSEM::IncludeThermoFluctuations]{}[Include\_Thermo\_Fluctuations][]

\begin{nmldescbody}
Specifies whether SEM fluctuations will be applied to the temperature and
density profiles based on the Strong Reynolds Analogy for temperature
fluctuations.
%
If this input is false, no fluctuations are added to the temperature and
density profiles.
%
See \sref{SEC:SEM:theory:ComputeBC} for details.
%
\end{nmldescbody}

%==========================================================================

\begin{nmldescitem}
Box_Num_Eddies_Multiplier = 1.0
\end{nmldescitem}
\inp*[nmlSEM::BoxNumEddiesMultiplier]{}[Box\_Num\_Eddies\_Multiplier][]

\begin{nmldescbody}
This input is used as a multiplier to the total number of synthetic eddies
used for the SEM boundary conditions, \thatis
%
%[{Number of eddies}] = [\inptext{Box\_Num\_Eddies\_Multiplier}] $\times$
%                       [{Volume of SEM box}] /
%                       [{Volume of synthetic eddy}].
%
\begin{gather*}
     [\text{Number of eddies}]
  =  [\inptext{Box\_Num\_Eddies\_Multiplier}] \times
     \frac{ [\text{Volume of SEM box}] }
          { [\text{Volume of synthetic eddy}] }
\end{gather*}
%
where the volumes are described in \sref{SEC:SEM:theory:Nomenclature}.
 
\NOTE 
The number of eddies for an individual SEM box are limited to be between a
minimum of 5,000 and a maximum of 500,000.
%
These limits were arbitrarily picked and hardwired into the SEM code, and
either limit can be modified in the source code if needed.
%
They were implemented to prevent
(1) too few eddies from being present to adequately model the flow, and
(2) too many eddies that add little to the modeling, but increase
    computational cost.
%
\end{nmldescbody}

%==========================================================================

\begin{nmldescitem}
Use_Restart_Eddies = T
\end{nmldescitem}
\inp*[nmlSEM::UseRestartEddies]{}[Use\_Restart\_Eddies][]

\begin{nmldescbody}
This input determines whether or not the SEM information contained in a restart
file is used or discarded after reading in the restart file.
%
It is recommended to use the SEM information in the restart file.
%
However, there are certain situations where starting with a fresh SEM is needed
or helpful which requires not using the SEM information in the restart file.
%
Probably the most common scenario is if the SEM convective velocity or length
scale is changed between simulations.
%
A change in either of these two values will result in a change to the size
of the SEM box, which will subsequently cause an abort error if any of the
restart eddies are no longer contained by the new SEM box.

\NOTE
Not using the SEM information in the restart file will result in a transient
pressure wave at the start of the restart simulation due to the discontinuity in
the SEM boundary condition between simulations.
\end{nmldescbody}

%==========================================================================

\FINALFIX{\newpage}
\begin{nmldescitem}
Tent_Function_Continuity = T
\end{nmldescitem}
\inp*[nmlSEM::TentFunctionContinuity]{}[Tent\_Function\_Continuity][]

\begin{nmldescbody}
This input specifies which tent function to use for the SEM.
%
Recall from \sref{SEC:SEM:theory:SumFluctuations} that the tent function
used by GFR is defined as
%
\begin{equation*}
  f_{\wedge}(a) = 
  \begin{cases}
    g(a) & \text{if } \abs{a} < 1 \\
    0    & \text{otherwise}
  \end{cases}
\end{equation*}
%
This input specifies the function $g(a)$ to use with the tent function, and the
input value generally corresponds to the continuity, $C^n$, of the overall tent
function:\\[-5pt]

\begin{tabular}{ @{} r @{ \qquad $g(a)=$} rl @{\quad : } l @{} }
  -1 & $\sqrt{\tfrac{     3}{    2}}$ & $\mspace{-17mu}\left(1-\abs{a}\right)$ & $C^0$ continuous (default) \\[5pt]
   0 & $\sqrt{\tfrac{    15}{   16}}$ & $\mspace{-17mu}\left(1-a^2\right)$     & $C^0$ continuous  \\[5pt]
   1 & $\sqrt{\tfrac{   315}{  256}}$ & $\mspace{-17mu}\left(1-a^2\right)^2$   & $C^1$ continuous  \\[5pt]
   2 & $\sqrt{\tfrac{  6006}{ 4096}}$ & $\mspace{-17mu}\left(1-a^2\right)^3$   & $C^2$ continuous  \\[5pt]
   3 & $\sqrt{\tfrac{109395}{65536}}$ & $\mspace{-17mu}\left(1-a^2\right)^4$   & $C^3$ continuous
\end{tabular}
\\[5pt]
\NOTE
The non-default tent functions have yet to be thoroughly tested.
\end{nmldescbody}

%==========================================================================

